%-------------------------------------------------------------------------------
% SECTION TITLE
%-------------------------------------------------------------------------------
\cvsection{Research and Education Experience}

%-------------------------------------------------------------------------------
% CONTENT
%-------------------------------------------------------------------------------
\begin{cventries}

%---------------------------------------------------------
  \cventry
    {Ph.D. Research Assistant | Advisor: \href{https://pitt-photonics.github.io/}{Nathan Youngblood}}
    {University of Pittsburgh}
    {Pittsburgh, PA}
    {Spring 2020 - present}
    {
      \begin{cvitems}
        \item {Developed an optical-injection-locking framework for distributed coherent photonic computing, combining semiconductor-laser rate-equation modeling, injection-ratio/detuning sweeps, coherent carrier recovery, and balanced coherent detection.}
        \item {Designed chip-scale photonic integrated circuits to translate fiber injection-locking systems onto integrated platforms, including waveguides, directional couplers, MMI/splitters, grating couplers, edge-coupling interfaces, and interferometric paths.}
        \item {Simulated silicon photonics and TFLN/LNOI passive and active photonic devices using GDSFactory, Cadence Virtuoso, KLayout, Lumerical, RSoft, and COMSOL workflows for layout, optical mode analysis, electro-optic modulation, and device optimization.}
        \item {Performed and supported photonic cleanroom fabrication and process evaluation, including spin coating/photoresist processing, e-beam lithography, etching, SEM inspection, and characterization of passive and active photonic structures.}
        \item {Integrated and aligned photonic chips using fiber arrays, grating-coupler alignment, edge coupling, optical packaging, VOAs, polarization controllers, and tunable-laser setups to connect fabricated devices to validation testbeds.}
        \item {Built high-speed EO/RF validation workflows using EOMs, high-speed photodetectors, balanced photodetection, OSAs, SVAs, spectrum analyzers, VNA-based characterization, RF detection, lock-in detection, PID phase feedback, and fiber-stretcher control.}
        \item {Automated photonic hardware validation with Python, PyVISA/SCPI, and NI MAX for wavelength, power, injection-ratio, and frequency sweeps; logged data and generated FFT/DSP analysis, plots, and repeatable reports.}
      \end{cvitems}
    }

%---------------------------------------------------------
  \cventry
    {Research Assistant | Advisor: \href{https://labsites.rochester.edu/zhang/}{Xi-Cheng Zhang}}
    {University of Rochester}
    {Rochester, NY}
    {Sep 2017 - Aug 2019}
    {
      \begin{cvitems}
        \item {Generated THz signals from liquid plasmas and quantified how temperature, polarity, and viscosity affected emission strength, bandwidth, repeatability, and spectral content.}
        \item {Processed time- and frequency-domain experimental data using SNR, sensitivity, regression, and repeatability analysis; prototyped digital cloaking with lenslet arrays and MATLAB pixel remapping.}
      \end{cvitems}
    }

%---------------------------------------------------------
  \cventry
    {Ph.D. in Electrical and Computer Engineering}
    {University of Pittsburgh}
    {Pittsburgh, PA}
    {Dec. 2026 expected}
    {}
\vspace{-0.4cm}

%---------------------------------------------------------
  \cventry
    {M.S. in Computer Science, Machine Learning track}
    {Georgia Institute of Technology}
    {Atlanta, GA}
    {May 2027 expected}
    {}
\vspace{-0.4cm}

%---------------------------------------------------------
  \cventry
    {M.S. in Optics (thesis track) \linebreak B.S. in Optical Engineering \linebreak B.S. in Applied Mathematics}
    {University of Rochester}
    {Rochester, NY}
    {Aug. 2019 \linebreak May 2017 \linebreak May 2017}
    {}
\vspace{-0.45cm}

\end{cventries}
